% =============================================================================
%  preamble.tex —— 全局样式
%  这里集中放宏包、字体、版式、定理环境、TikZ 样式、代码样式。
%  正文请写在各 sections/*.tex 里；改外观只改这个文件即可。
%  编译： latexmk -xelatex main.tex   （见 Makefile / build.sh）
% =============================================================================


% -----------------------------------------------------------------------------
% 1. 页面与行距
% -----------------------------------------------------------------------------
\usepackage{geometry}
\geometry{
  a4paper,
  left   = 3.2cm, right  = 3.2cm,
  top    = 3.0cm, bottom = 2.8cm,
  headsep = 0.6cm, footskip = 1.0cm
}
\usepackage{setspace}
\setstretch{1.42}          % 中文行距：1.4～1.5 之间比较舒服，可按需调

% 中文段落断行：给 TeX 留一点额外弹性，减少长公式/长词造成的 overfull
\setlength{\emergencystretch}{2.5em}
\tolerance         = 2000
\hbadness          = 10000   % 只关心 overfull，忽略轻微的 underfull
\vbadness          = 10000


% -----------------------------------------------------------------------------
% 2. 字体（XeLaTeX + fontspec）
%    西文用 Latin Modern，与默认的 Computer Modern 数学字体完全统一；
%    中文正文宋体、标题/图表题黑体，与参考笔记的观感一致。
%    若想换成 Times，把 \setmainfont 换成 Times New Roman 即可。
% -----------------------------------------------------------------------------
\usepackage{fontspec}
% 西文衬线族（正文）
\setmainfont{lmroman10-regular.otf}[
  BoldFont       = lmroman10-bold.otf,
  ItalicFont     = lmroman10-italic.otf,
  BoldItalicFont = lmroman10-bolditalic.otf,
  Ligatures      = TeX
]
% 西文无衬线族（标题、图表题）
\setsansfont{Helvetica Neue}
% 西文等宽族（代码）
\setmonofont{Menlo}[Scale=0.92]

% 中文字体
% 注：ctex 已经定义了 \songti \heiti \kaishu \fangsong 等切换命令，
%     这里只需要把四套字体族指向 macOS 上对应的字体即可。
%     Heiti SC 带有真正的 Bold 字面，标题不会出现缺字警告。
\setCJKmainfont{Songti SC}[ItalicFont=STFangsong]   % 宋体；强调用仿宋
\setCJKsansfont{Heiti SC}                          % 黑体（标题）
\setCJKmonofont{STFangsong}                        % 仿宋（等宽场景）


% -----------------------------------------------------------------------------
% 3. 数学
% -----------------------------------------------------------------------------
\usepackage{amsmath}
\usepackage{amssymb}
\usepackage{mathtools}
\usepackage{bm}
\allowdisplaybreaks

% 常用的数学记号（避免正文里反复写长命令）
\providecommand{\E}{\mathbb{E}}                        % 期望
\newcommand{\Ncal}{\mathcal{N}}                        % 高斯分布
\newcommand{\Dkl}{\mathrm{D}_{\mathrm{KL}}}            % KL 散度
\newcommand{\bx}{\bm{x}}
\newcommand{\bz}{\bm{z}}
\newcommand{\beps}{\bm{\varepsilon}}
\newcommand{\bI}{\bm{I}}
\newcommand{\bzero}{\bm{0}}
\newcommand{\bmu}{\bm{\mu}}
\newcommand{\abt}{\bar{\alpha}_t}                      % \bar\alpha_t
\newcommand{\abs}{\bar{\alpha}_s}
\DeclareMathOperator*{\argmax}{arg\,max}
\DeclareMathOperator*{\argmin}{arg\,min}


% -----------------------------------------------------------------------------
% 4. 颜色（先定义，后面 TikZ / 表格 / 代码都要用）
% -----------------------------------------------------------------------------
\usepackage{xcolor}
\definecolor{dmNavy} {HTML}{2B2D5C}   % 深藏青：主线条、标题
\definecolor{dmSteel}{HTML}{3C6E9F}   % 钢蓝：强调、链接
\definecolor{dmLav}  {HTML}{DFDBF4}   % 淡紫：节点底色（对齐参考笔记）
\definecolor{dmSky}  {HTML}{D6E4F7}   % 淡蓝
\definecolor{dmMint} {HTML}{D6EDE7}   % 淡青
\definecolor{dmRose} {HTML}{F7DFDF}   % 淡粉
\definecolor{dmGray} {HTML}{6B6F8A}   % 中性灰


% -----------------------------------------------------------------------------
% 5. 版式细节：标题居中、页码在右上角、图表题样式
% -----------------------------------------------------------------------------
\usepackage{fancyhdr}
\pagestyle{fancy}
\fancyhf{}
\renewcommand{\headrulewidth}{0pt}
\fancyhead[R]{\small\thepage}

% 章节标题：一级居中加粗黑体，二级左对齐
\ctexset{
  section = {
    format     = \sffamily\bfseries\zihao{3}\centering,
    aftername  = \quad,
    beforeskip = 2.0ex plus .3ex minus .2ex,
    afterskip  = 1.4ex plus .2ex,
  },
  subsection = {
    format     = \sffamily\bfseries\zihao{4},
    aftername  = \quad,
    beforeskip = 2.0ex plus .3ex minus .2ex,
    afterskip  = 1.0ex plus .2ex,
  },
  subsubsection = {
    format     = \sffamily\bfseries\zihao{5},
    aftername  = \quad,
    beforeskip = 1.6ex plus .3ex minus .2ex,
    afterskip  = 0.8ex plus .2ex,
  },
}

% 浮动体、图表题
\usepackage{graphicx}
\usepackage{float}
\usepackage{booktabs}
\usepackage{array}
\usepackage{tabularx}
\usepackage{multirow}
\usepackage{caption}
\captionsetup{
  font            = {sf, bf, small},
  labelsep        = colon,
  justification   = centering,
  singlelinecheck = true,
  skip            = 6pt,
}
\captionsetup[table]{position=top}
\newcommand{\tabref}[1]{表~\ref{#1}}
\newcommand{\figref}[1]{图~\ref{#1}}

% 列表间距
\usepackage{enumitem}
\setlist[itemize]{itemsep=0.25ex plus .1ex, topsep=0.6ex plus .2ex, parsep=0pt}
\setlist[enumerate]{itemsep=0.25ex plus .1ex, topsep=0.6ex plus .2ex, parsep=0pt}


% -----------------------------------------------------------------------------
% 6. 定理类环境（正文一律正体，中文不用斜体）
% -----------------------------------------------------------------------------
\usepackage{amsthm}
\newtheoremstyle{cn}
  {1.3ex plus .2ex}% 上方间距
  {1.3ex plus .2ex}% 下方间距
  {\normalfont}%    正文字体：正体
  {}%              缩进
  {\sffamily\bfseries}% 头部字体
  {}%              头部标点
  {0.7em}%         头部后间距
  {}%              头部规格
\theoremstyle{cn}
\newtheorem{definition} {定义}[section]
\newtheorem{proposition}[definition]{命题}
\newtheorem{theorem}    [definition]{定理}
\newtheorem{corollary}  [definition]{推论}
\newtheorem{example}    [definition]{例}
\newtheorem{remark}     [definition]{注}

% 自带的“证明”环境（比 amsthm 的 proof 更适合中文排版）
\newenvironment{pf}[1][证明]{%
  \par\vspace{0.3ex}%
  \noindent{\sffamily\bfseries #1}\quad\ignorespaces
}{%
  \unskip\nobreak\hfill$\square$\par\vspace{1.0ex}%
}


% -----------------------------------------------------------------------------
% 7. TikZ / pgfplots
% -----------------------------------------------------------------------------
\usepackage{tikz}
\usepackage{pgfplots}
\pgfplotsset{compat=1.18}
\usetikzlibrary{
  arrows.meta, positioning, calc, fit, backgrounds,
  shapes.geometric, shapes.misc, decorations.pathreplacing
}

\tikzset{
  % 通用节点：圆角、描边 + 淡紫底（和参考笔记的示意图风格一致）
  dbox/.style  = {draw=dmNavy, line width=0.7pt, fill=dmLav,
                  rounded corners=3pt, inner sep=4pt, minimum height=2.0em,
                  font=\small\sffamily, text=dmNavy, align=center},
  dbox2/.style = {dbox, fill=dmSky},
  dbox3/.style = {dbox, fill=dmMint},
  dbox4/.style = {dbox, fill=dmRose},
  % 打通的箭头
  darr/.style  = {-{Latex[length=2.1mm, width=1.5mm]}, draw=dmNavy, line width=0.7pt},
  dbarr/.style = {{Latex[length=2.1mm, width=1.5mm]}-{Latex[length=2.1mm, width=1.5mm]},
                  draw=dmNavy, line width=0.7pt},
  % 注释文字
  dlab/.style  = {font=\footnotesize\sffamily, text=dmNavy, inner sep=2pt},
  dnum/.style  = {font=\small\itshape, text=dmNavy},
}


% -----------------------------------------------------------------------------
% 8. 算法
% -----------------------------------------------------------------------------
\usepackage{algorithm}
\usepackage{algpseudocode}
\floatname{algorithm}{算法}
\renewcommand{\algorithmicrequire}{\textbf{输入：}}
\renewcommand{\algorithmicensure}{\textbf{输出：}}
\renewcommand{\algorithmicforall}{\textbf{对}}
\renewcommand{\algorithmicdo}{\textbf{do}}
\renewcommand{\algorithmicend}{\textbf{end}}
% 让算法环境的说明文字也有统一的字号
\makeatletter
\renewcommand{\ALG@beginalgorithmic}{\small}
\makeatother


% -----------------------------------------------------------------------------
% 9. 强调框 / 代码
% -----------------------------------------------------------------------------
\usepackage[most]{tcolorbox}
\newtcolorbox{keybox}[1][要点]{
  enhanced, breakable,
  colback  = dmLav!30, colframe = dmNavy,
  boxrule  = 0.7pt, arc = 3pt,
  left = 7pt, right = 7pt, top = 6pt, bottom = 6pt,
  fonttitle = \sffamily\bfseries,
  coltitle  = dmNavy, colbacktitle = dmLav!70,
  title = {#1},
  attach boxed title to top left = {xshift=8pt, yshift=-2pt},
  boxed title style = {boxrule=0pt, arc=2pt},
  top = 12pt,
}

\usepackage{listings}
\lstdefinestyle{pystyle}{
  language        = Python,
  basicstyle      = \ttfamily\footnotesize,
  keywordstyle    = \color{dmSteel}\bfseries,
  commentstyle    = \color{dmGray}\itshape,
  stringstyle     = \color{dmNavy},
  numbers         = left,
  numberstyle     = \tiny\color{dmGray},
  numbersep       = 8pt,
  frame           = single,
  rulecolor       = \color{dmGray!45},
  backgroundcolor = \color{dmGray!6},
  breaklines      = true,
  showstringspaces= false,
  tabsize         = 4,
  xleftmargin     = 1.6em,
  framexleftmargin= 1.6em,
  columns         = fullflexible,
  keepspaces      = true,
  upquote         = true,
}
\lstset{style=pystyle}


% -----------------------------------------------------------------------------
% 10. 超链接与参考文献
%     （hyperref 要放在最后；biblatex 放在 hyperref 之后）
% -----------------------------------------------------------------------------
\usepackage{hyperref}
\hypersetup{
  colorlinks  = true,
  linkcolor   = dmNavy,
  citecolor   = dmSteel,
  urlcolor    = dmSteel,
  bookmarksnumbered = true,
  pdftitle    = {扩散模型基础},
  pdfsubject  = {Diffusion Models, DDPM, 变分推断},
}

\usepackage[
  backend  = biber,
  style    = numeric-comp,
  sortcites= true,
  giveninits = true,
  maxbibnames = 6,
  doi = false, isbn = false,
  url = true
]{biblatex}
\addbibresource{refs.bib}
\setlength{\bibitemsep}{0.4ex}
